\begin{abstract}
Visual question answering (VQA) systems for blind and low-vision users must do more than return an answer: they must also know when an answer is unreliable, explain why the input failed, and help the user recover. This paper studies a lightweight reliability layer for assistive VQA that combines answerability triage, human-labeled image-quality defect diagnosis, selective-prediction diagnostics, and actionable retake guidance. Using VizWiz-VQA and VizWiz-QualityIssues, we join real answerability labels with quality defects including blur, brightness, darkness, obstruction, framing, rotation, and unrecognizable content. Frozen visual embeddings from \clip, \dinov, and \mobilenet are used to train small answerability and defect heads, while a frozen \vilt VQA model provides answer confidence for selective-prediction analysis.

The strongest triage head reaches AUROC $0.7926$ with \clip; the strongest defect diagnosis result reaches mAP $0.6309$ with \dinov; and the strongest actionable recovery result reaches \arr{} $0.7196$ with \clip. However, defect-aware risk scores do not significantly improve area under the risk-coverage curve over the VQA model's own confidence, and even oracle ground-truth defect features are worse than global confidence in this setup. We therefore argue that visual defect diagnosis is best viewed not as a replacement for confidence-based risk ranking, but as an interpretable refusal and recovery mechanism: confidence indicates whether to trust an answer, while defect diagnosis explains why the system should refuse and what the user can do next.
\end{abstract}
